% GENERATED FILE. Edit canonical lessons, then run npm run generate:books.
% canonical-source-hash: fnv1a64:45298c28b08a2f57
% canonical-lessons: GU-C244-ganit, GU-C244-bhugol, GU-C244-nibandh, GU-C244-samjavvu, GU-C244-ukelvu

\chapter{Maths, maps and essays}
\label{ch:gu-maths-maps-and-essays}
\hlchaptermodality{244}

\begin{chapteropening}
\textbf{By the end of this chapter:} \emph{I can name school subjects and say who explains and solves problems.}

Complete the last lesson of chapter 244: I can name school subjects and say who explains and solves problems.
\end{chapteropening}

\section[\texorpdfstring{\gu{ગણિત} (gaṇit)}{ગણિત (gaṇit)}]{\gu{ગણિત} (gaṇit) — mathematics, maths}

\label{lesson:GU-C244-ganit}



\begin{quote}
Before the new one: say the Gujarati for school bag, then the Gujarati for present.
\end{quote}

\subsection*{You'll want to know: \gu{ગણિત}}
\textbf{\gu{ગણિત}} — \emph{gaṇit} — \textquotedblleft{}mathematics, maths\textquotedblright{}.

\begin{quote}
\textbf{\gu{મને} \gu{ગણિત} \gu{ગમે} \gu{છે}.} — \emph{mane gaṇit game che} — \textquotedblleft{}I like maths.\textquotedblright{}
\end{quote}

A neuter noun.

\begin{cousinweb}
\textbf{\gu{ગણિત}} (\emph{gaṇit}) is Sanskrit \emph{gaṇita}, \textquotedblleft{}counted, reckoned\textquotedblright{}, taken over whole as a learned loan, a tatsama (CDIAL 3996). The taught verb \emph{gaṇvũ}, \textquotedblleft{}to count\textquotedblright{}, is built on the same root.
\end{cousinweb}

\begin{grammarlens}[title={What you've built}]
One more word for everyday life at A2.
\end{grammarlens}

\subsection*{Your turn}
\begin{itemize}
\raggedright
  \item \emph{Say it:} \emph{gaṇit}
  \item \emph{Say it:} \emph{gaṇit}, once more
  \item \emph{Say it:} \emph{mane gaṇit game che} — I like maths.
  \item \emph{Recall:} say the Gujarati for school bag, then the Gujarati for present, then say \emph{gaṇit} again
\end{itemize}

\begin{tcolorbox}[breakable,colback=teal!4,colframe=teal!35!black,title={Before you move on}]
What does \gu{ગણિત} mean? (\textquotedblleft{}Mathematics, maths\textquotedblright{}.) Say it once more.
\end{tcolorbox}
\section[\texorpdfstring{\gu{ભૂગોળ} (bhūgoḷ)}{ભૂગોળ (bhūgoḷ)}]{\gu{ભૂગોળ} (bhūgoḷ) — geography}

\label{lesson:GU-C244-bhugol}



\begin{quote}
Before the new one: say the Gujarati for present, then the Gujarati for mathematics.
\end{quote}

\subsection*{You'll want to know: \gu{ભૂગોળ}}
\textbf{\gu{ભૂગોળ}} — \emph{bhūgoḷ} — \textquotedblleft{}geography\textquotedblright{}.

\begin{quote}
\textbf{\gu{અમે} \gu{ભૂગોળમાં} \gu{નકશા} \gu{ભણીએ} \gu{છીએ}.} — \emph{ame bhūgoḷmā̃ nakśā bhaṇīe chīe} — \textquotedblleft{}We study maps in geography.\textquotedblright{}
\end{quote}

A feminine noun.

\begin{cousinweb}
\textbf{\gu{ભૂગોળ}} (\emph{bhūgoḷ}) is a learned loan from Sanskrit \emph{bhūgola}, \textquotedblleft{}the globe\textquotedblright{}, which EtymDB lists as borrowed from Sanskrit. It is built from \emph{bhū}, \textquotedblleft{}earth\textquotedblright{}, and \emph{gola}, \textquotedblleft{}ball\textquotedblright{} (CDIAL 4321).
\end{cousinweb}

\begin{grammarlens}[title={What you've built}]
One more word for everyday life at A2.
\end{grammarlens}

\subsection*{Your turn}
\begin{itemize}
\raggedright
  \item \emph{Say it:} \emph{bhūgoḷ}
  \item \emph{Say it:} \emph{bhūgoḷ}, once more
  \item \emph{Say it:} \emph{ame bhūgoḷmā̃ nakśā bhaṇīe chīe} — We study maps in geography.
  \item \emph{Recall:} say the Gujarati for present, then the Gujarati for mathematics, then say \emph{bhūgoḷ} again
\end{itemize}

\begin{tcolorbox}[breakable,colback=teal!4,colframe=teal!35!black,title={Before you move on}]
What does \gu{ભૂગોળ} mean? (\textquotedblleft{}Geography\textquotedblright{}.) Say it once more.
\end{tcolorbox}
\section[\texorpdfstring{\gu{નિબંધ} (nibandh)}{નિબંધ (nibandh)}]{\gu{નિબંધ} (nibandh) — essay}

\label{lesson:GU-C244-nibandh}



\begin{quote}
Before the new one: say the Gujarati for mathematics, then the Gujarati for geography.
\end{quote}

\subsection*{You'll want to know: \gu{નિબંધ}}
\textbf{\gu{નિબંધ}} — \emph{nibandh} — \textquotedblleft{}essay\textquotedblright{}.

\begin{quote}
\textbf{\gu{તે} \gu{ગુજરાતીમાં} \gu{નિબંધ} \gu{લખે} \gu{છે}.} — \emph{te gujarātīmā̃ nibandh lakhe che} — \textquotedblleft{}She writes an essay in Gujarati.\textquotedblright{}
\end{quote}

A masculine noun.

\begin{cousinweb}
\textbf{\gu{નિબંધ}} (\emph{nibandh}) is Sanskrit \emph{nibandha}, \textquotedblleft{}fastening\textquotedblright{} and later \textquotedblleft{}literary composition\textquotedblright{}, taken over whole as a tatsama (CDIAL 7222).
\end{cousinweb}

\begin{grammarlens}[title={What you've built}]
One more word for everyday life at A2.
\end{grammarlens}

\subsection*{Your turn}
\begin{itemize}
\raggedright
  \item \emph{Say it:} \emph{nibandh}
  \item \emph{Say it:} \emph{nibandh}, once more
  \item \emph{Say it:} \emph{te gujarātīmā̃ nibandh lakhe che} — She writes an essay in Gujarati.
  \item \emph{Recall:} say the Gujarati for mathematics, then the Gujarati for geography, then say \emph{nibandh} again
\end{itemize}

\begin{tcolorbox}[breakable,colback=teal!4,colframe=teal!35!black,title={Before you move on}]
What does \gu{નિબંધ} mean? (\textquotedblleft{}Essay\textquotedblright{}.) Say it once more.
\end{tcolorbox}
\section[\texorpdfstring{\gu{સમજાવવું} (samjāvvũ)}{સમજાવવું (samjāvvũ)}]{\gu{સમજાવવું} (samjāvvũ) — to explain}

\label{lesson:GU-C244-samjavvu}



\begin{quote}
Before the new one: say the Gujarati for geography, then the Gujarati for essay.
\end{quote}

\subsection*{You'll want to know: \gu{સમજાવવું}}
\textbf{\gu{સમજાવવું}} — \emph{samjāvvũ} — \textquotedblleft{}to explain\textquotedblright{}.

\begin{quote}
\textbf{\gu{ગુરુ} \gu{ગણિત} \gu{સમજાવે} \gu{છે}.} — \emph{guru gaṇit samjāve che} — \textquotedblleft{}The teacher explains maths.\textquotedblright{}
\end{quote}

The causative of \textbf{\gu{સમજવું}}, \emph{samajvũ}, \textquotedblleft{}to understand\textquotedblright{}: to make someone understand.

\begin{cousinweb}
\textbf{\gu{સમજાવવું}} (\emph{samjāvvũ}) is built on the taught \emph{samajvũ}, \textquotedblleft{}to understand\textquotedblright{}, with the causative \emph{-āv-}. CDIAL derives \emph{samajvũ} from Sanskrit \emph{saṁbudhyatē}, \textquotedblleft{}wakes up, understands\textquotedblright{} (CDIAL 12959).
\end{cousinweb}

\begin{grammarlens}[title={What you've built}]
One more word for everyday life at A2.
\end{grammarlens}

\subsection*{Your turn}
\begin{itemize}
\raggedright
  \item \emph{Say it:} \emph{samjāvvũ}
  \item \emph{Say it:} \emph{samjāvvũ}, once more
  \item \emph{Say it:} \emph{guru gaṇit samjāve che} — The teacher explains maths.
  \item \emph{Recall:} say the Gujarati for geography, then the Gujarati for essay, then say \emph{samjāvvũ} again
\end{itemize}

\begin{tcolorbox}[breakable,colback=teal!4,colframe=teal!35!black,title={Before you move on}]
What does \gu{સમજાવવું} mean? (\textquotedblleft{}To explain\textquotedblright{}.) Say it once more.
\end{tcolorbox}
\section[\texorpdfstring{\gu{ઉકેલવું} (ukelvũ)}{ઉકેલવું (ukelvũ)}]{\gu{ઉકેલવું} (ukelvũ) — to solve, to work out}

\label{lesson:GU-C244-ukelvu}



\begin{quote}
Before the new one: say the Gujarati for essay, then the Gujarati for to explain.
\end{quote}

\subsection*{You'll want to know: \gu{ઉકેલવું}}
\textbf{\gu{ઉકેલવું}} — \emph{ukelvũ} — \textquotedblleft{}to solve, to work out\textquotedblright{}.

\begin{quote}
\textbf{\gu{તે} \gu{અઘરો} \gu{પ્રશ્ન} \gu{ઉકેલે} \gu{છે}.} — \emph{te agharo praśna ukele che} — \textquotedblleft{}He solves a hard problem.\textquotedblright{}
\end{quote}

\begin{cousinweb}
\textbf{\gu{ઉકેલવું}} (\emph{ukelvũ}) is a tadbhava. It comes from \emph{*utkēlyatē}, \textquotedblleft{}is disentangled\textquotedblright{}, a reconstructed form; CDIAL glosses the Gujarati verb \textquotedblleft{}to disentangle, make out a badly written letter\textquotedblright{} (CDIAL 1734).
\end{cousinweb}

\begin{grammarlens}[title={What you've built}]
One more word for everyday life at A2.
\end{grammarlens}

\subsection*{Your turn}
\begin{itemize}
\raggedright
  \item \emph{Say it:} \emph{ukelvũ}
  \item \emph{Say it:} \emph{ukelvũ}, once more
  \item \emph{Say it:} \emph{te agharo praśna ukele che} — He solves a hard problem.
  \item \emph{Recall:} say the Gujarati for essay, then the Gujarati for to explain, then say \emph{ukelvũ} again
\end{itemize}

\begin{tcolorbox}[breakable,colback=teal!4,colframe=teal!35!black,title={Before you move on}]
What does \gu{ઉકેલવું} mean? (\textquotedblleft{}To solve, to work out\textquotedblright{}.) Say it once more.
\end{tcolorbox}
